\documentclass{article}

\PassOptionsToPackage{numbers}{natbib}
\usepackage[preprint]{neurips_2026}

\usepackage[utf8]{inputenc}
\usepackage[T1]{fontenc}
\usepackage{hyperref}
\usepackage{url}
\usepackage{booktabs}
\usepackage{amsfonts}
\usepackage{amsmath}
\usepackage{amssymb}
\usepackage{nicefrac}
\usepackage{microtype}
\usepackage{xcolor}
\usepackage{graphicx}
\usepackage{enumitem}
\usepackage{multirow}
\usepackage{array}
\usepackage{longtable}

\definecolor{good}{RGB}{0,110,60}
\definecolor{warn}{RGB}{175,105,0}
\definecolor{bad}{RGB}{175,0,0}

\newcommand{\dok}[1]{\textcolor{good}{#1}}
\newcommand{\dwarn}[1]{\textcolor{warn}{#1}}
\newcommand{\dbad}[1]{\textcolor{bad}{#1}}
\newcommand{\pu}{\,\text{pu}}

\title{The \texttt{\_ppNR\_} Dataset Family, v2\\
\large 31 grids under one configuration: clean DC start, perturbed static
generation, and a PV-bus reactive entry consistent with the stored voltage}

\author{Changhun Kim}
\date{August 16, 2026}

\begin{document}
\maketitle

\begin{abstract}
This document describes the \textbf{v2} \texttt{\_ppNR\_} dataset family: 31
parquet files covering 31 power grids from 4 to 9{,}241 buses, every one
generated under an \emph{identical} configuration for the first time.  The v1
family documented previously was internally inconsistent in three ways, each of
which is measured and corrected here.  (i)~The 28 PPC grids carried
$5^\circ/0.02\pu$ noise on \texttt{u\_start} while the 3 CGMES grids did not;
the noise is now off everywhere, because it degrades \texttt{u\_start} as a
predictor of \texttt{u\_newton} by $35$--$134\%$ while adding no information.
(ii)~\texttt{net.sgen} was never perturbed, freezing $56.2\%$ of scheduled
generation on case24\_ieee\_rts, $34.5\%$ on case89pegase, $17.2\%$ on LVN heo1
and $100\%$ on SimBench across all scenarios; static generation now follows the
load scale and the dispatch jitter.  (iii)~The PV-bus reactive entry of
\texttt{S\_start} was back-computed at a flat start voltage that is \emph{not}
the \texttt{u\_start} stored beside it, inconsistently so per grid; it is now
recomputed from \texttt{u\_start}.  We also document the per-grid load-scale
feasibility cliff that explains every shortfall yield in v1 to within one
percentage point, and a systematic DC-versus-AC angle bias that accumulates
with distance from the slack.
\end{abstract}


% ===================================================================
\section{Location and Selection}
\label{sec:sel}
% ===================================================================

All files reside in a single directory on the FAU Fritz cluster:

\begin{center}
\texttt{/home/vault/iwi5/iwi5295h/PIGNN-Attn-LS/ScenarioSynthesis\_PPC/out/}
\end{center}

\noindent
\textbf{The v1 selection rule no longer works.}  Selecting on the
\texttt{\_ppNR\_} solver tag alone now matches several generations of the same
grids.  The v2 family is selected by its configuration tag:

\begin{center}
\texttt{ls .../out/*\_u0clean\_sgenpert\_pvqvstart\_*directSI.parquet}
\end{center}

\noindent
This yields exactly 31 files.  The tag is authoritative: it is emitted by the
filename builder only when all three of \texttt{--rand\_u\_start} (absent),
\texttt{--perturb\_sgen} and \texttt{--pv\_q\_from\_vstart} hold, so no earlier
generation can be selected by accident.  Anything matching \texttt{\_cNR\_} is a
different family entirely --- it was labelled with the in-repo custom solver,
not pandapower --- and is excluded automatically.


% ===================================================================
\section{Filename Scheme}
% ===================================================================

Filenames encode the full experimental configuration so that runs differing in
any controlled variable cannot collide:

\begin{center}
\small
\texttt{<grid>\_<ybus>\_<preset>\_<start>\_<solver>\_<load>\_}\\
\texttt{<u0>[\_<flags>]\_<unit>\_<rows>\_NR\_branchrows\_directSI.parquet}
\end{center}

\noindent
For example:

\begin{center}
\scriptsize
\texttt{case118\_ppcY\_backbone\_dc\_compile\_ppNR\_ls0.60-1.40\_}\\
\texttt{u0clean\_sgenpert\_pvqvstart\_siNR\_37000\_NR\_branchrows\_directSI.parquet}
\end{center}

\begin{center}
\begin{tabular}{ll}
\toprule
\textbf{Field} & \textbf{Meaning} \\
\midrule
\texttt{ppcY}            & Y-bus taken from the pandapower PPC (not re-stamped) \\
\texttt{backbone}        & scenario preset (Section~\ref{sec:pert}) \\
\texttt{dc\_compile}     & initial voltage from a DC power-flow solve \\
\texttt{ppNR}            & label solver = pandapower Newton--Raphson \\
\texttt{ls0.60-1.40}     & effective global load-scale range ($\pm40\%$) \\
\texttt{u0clean}         & \textbf{no start-point jitter} (Section~\ref{sec:u0}) \\
\texttt{sgenpert}        & \texttt{net.sgen} is perturbed (Section~\ref{sec:sgen}) \\
\texttt{pvqvstart}       & PV-bus $Q$ from \texttt{u\_start} (Section~\ref{sec:pvq}) \\
\texttt{siNR}            & NR solved and stored in SI units (not per-unit) \\
\texttt{37000}           & \textbf{true row count} in the file \\
\bottomrule
\end{tabular}
\end{center}

\paragraph{The row-count field is now trustworthy.}
In v1 it was built from \texttt{--runs}, which counts \emph{attempts}.  With
\texttt{--drop\_nonconverged} the two differ, so
\texttt{case145\_...\_36000\_...} actually held $7{,}885$ rows and
\texttt{case300\_...\_36000\_...} held $27{,}994$.  Merged array datasets were
named by the merger and did carry true counts, so the two halves of the v1
corpus meant different things by the same field.  The driver now renames the
file to the true retained count on close, and this is verified against the
parquet metadata for all 31 files.


% ===================================================================
\section{Perturbation: the \texttt{backbone} Preset}
\label{sec:pert}
% ===================================================================

Every file uses \texttt{--scenario\_level backbone} with no CLI overrides of the
load range.  This is a \emph{fixed-Y-bus} perturbation family: topology and
admittance are untouched, so all rows of a grid share one admittance matrix and
the dataset remains \texttt{share\_grid}-compatible.

\begin{center}
\begin{tabular}{lcl}
\toprule
\textbf{Quantity} & \textbf{Value} & \textbf{Effect} \\
\midrule
Global load scale        & $[0.60,\ 1.40]$      & light $\to$ heavy $\to$ near-collapse \\
Scale gen.\ with load    & yes                  & keeps slack burden bounded \\
\textbf{Scale sgen with load} & \textbf{yes}    & \textbf{new in v2} (Section~\ref{sec:sgen}) \\
Per-bus load $P$ jitter  & $\sigma_P = 0.10$    & decorrelates buses \\
Per-bus load $Q$ jitter  & $\sigma_Q = 0.15$    & reactive less predictable than active \\
Generator jitter         & $\sigma_G = 0.05$    & dispatch variation, \texttt{gen} \emph{and} \texttt{sgen} \\
PV voltage setpoint      & $[0.95,\ 1.05]$\,pu  & realistic setpoint band \\
Start-point jitter       & \textbf{off}         & \textbf{changed in v2} (Section~\ref{sec:u0}) \\
\midrule
Line outage probability  & $0.0$                & \textbf{topology fixed} \\
Admittance jitter        & $0.0$                & \textbf{Y-bus fixed} \\
\bottomrule
\end{tabular}
\end{center}

\noindent
Note the ordering, which matters for Section~\ref{sec:u0}: every quantity above
is applied to the network \emph{before} the DC power flow that produces
\texttt{u\_start}.  They perturb the \emph{scenario}, and \texttt{u\_start} and
\texttt{u\_newton} move together in response.

\paragraph{Correction to the v1 report.}
The v1 document asserted that ``narrowing [the load range] to $\pm20\%$ or
$\pm10\%$ \dots changed convergence by less than $6\%$, establishing that
solver choice rather than perturbation magnitude governs yield.''  That
measurement was made with the custom solver.  Under pandapower NR the load
range governs yield almost completely, as Section~\ref{sec:cliff} shows.


% ===================================================================
\section{What Changed in v2, and Why}
% ===================================================================

\subsection{The start-point jitter is off}
\label{sec:u0}

In v1 the 28 PPC files were generated while \texttt{rand\_u\_start} was read
from the preset, which sets it on ($5^\circ$ angle jitter on all non-slack
buses, $0.02\pu$ magnitude jitter on PQ buses).  The generator was later
changed to make \texttt{rand\_u\_start} CLI-only and default off, so the three
CGMES files, generated afterwards, had a clean DC start.  The v1 report flagged
this as ``the one substantive inconsistency in the family.''

\paragraph{It is not merely an inconsistency; the jitter is harmful.}
The learning task is $\texttt{u\_start} \mapsto \texttt{u\_newton}$, and
$\texttt{u\_newton}$ is a function of $(S, Y)$ alone --- it does not depend on
the initial guess.  Since the model already receives $S$ and $Y$,
\texttt{u\_start} carries no information about the target.  Its entire value is
being a \emph{cheap approximation of the answer}, which makes corrupting it a
pure loss.  Measured mean angle error of \texttt{u\_start} against
\texttt{u\_newton}, with the per-grid bias of Section~\ref{sec:dcbias} removed:

\begin{center}
\begin{tabular}{lrrr}
\toprule
\textbf{Grid} & \textbf{clean} & \textbf{jittered} & \textbf{degradation} \\
\midrule
case118        & $1.145^\circ$ & $2.675^\circ$ & \dbad{$+134\%$} \\
case300        & $1.959^\circ$ & $3.216^\circ$ & \dbad{$+64\%$} \\
case1354pegase & $2.997^\circ$ & $4.041^\circ$ & \dbad{$+35\%$} \\
\bottomrule
\end{tabular}
\end{center}

\noindent
The jitter is also the wrong instrument for start-point diversity.  A DC power
flow is a linear solve, $\theta = B^{-1}P$; it takes no initial guess, so noise
applied \emph{before} it is simply overwritten.  What does reach the DC
solution is the scenario perturbation of Section~\ref{sec:pert} --- and PV-bus
magnitudes already carry it, since the DC start sets $|V| = 1.0\pu$ everywhere
and then overwrites generator buses with their (jittered) setpoints.  Only PQ
magnitudes were flat, so the $\pm2\%$ magnitude jitter was adding noise to a
constant, and the $\pm5^\circ$ angle jitter was corrupting the one genuinely
informative component.

\paragraph{Verification.}
Because a clean DC start leaves every PQ bus at exactly $1.0\pu$, the absence of
jitter is an exact check rather than a threshold.  Across all 31 files the
maximum deviation of $|u_{\text{start}}|$ from $1.0\pu$ at PQ buses is
$\sim\!10^{-16}$, i.e.\ machine epsilon.

\paragraph{An A/B pair is retained.}
\texttt{case118} and \texttt{case300} exist in both arms --- clean and jittered,
identical in every other respect --- so the claim above can be tested on a
trained model rather than argued.  These four files are kept outside the v2
family and are not selected by the tag in Section~\ref{sec:sel}.


\subsection{Static generation is perturbed}
\label{sec:sgen}

\texttt{net.sgen} holds static generators: wind, PV plants, CHP --- PQ
injections with a specified $p$ and $q$.  In v1 the perturbation code never
referenced \texttt{net.sgen} at all.  The global load scale and the dispatch
jitter reached \texttt{net.gen} and \texttt{net.load} only, so \emph{static
generation was identical in all $36{,}000$ scenarios of every grid}.

This is negligible where sgen is negligible and decisive where it is not:

\begin{center}
\begin{tabular}{lrrr}
\toprule
\textbf{Grid} & \textbf{sgen units} & \textbf{sgen MW} & \textbf{share of scheduled generation} \\
\midrule
SimBench          & 121   & $3{,}715.6$  & \dbad{$100.0\%$} \\
case24\_ieee\_rts & 22    & $1{,}632.2$  & \dbad{$56.2\%$} \\
case89pegase      & 6     & $2{,}430.8$  & \dbad{$34.5\%$} \\
case3120sp        & 257   & $5{,}235.7$  & \dbad{$25.1\%$} \\
case6495rte       & 1{,}157 & $19{,}139.4$ & \dwarn{$17.9\%$} \\
case6470rte       & 1{,}125 & $19{,}428.9$ & \dwarn{$17.5\%$} \\
case6515rte       & 1{,}158 & $20{,}810.2$ & \dwarn{$17.5\%$} \\
LVN heo1          & 33    & $36.9$       & \dwarn{$17.2\%$} \\
case5             & 1     & $170.0$      & \dwarn{$17.0\%$} \\
case9241pegase    & 434   & $23{,}055.8$ & \dwarn{$6.8\%$} \\
case2848rte       & 300   & $2{,}743.1$  & $5.0\%$ \\
case2869pegase    & 180   & $6{,}497.6$  & $4.7\%$ \\
case1354pegase    & 52    & $1{,}086.3$  & $1.5\%$ \\
case300           & 8     & $321.8$      & $1.4\%$ \\
case1888rte       & 82    & $679.9$      & $1.1\%$ \\
\bottomrule
\end{tabular}
\end{center}

\noindent
Thirteen of 28 PPC grids and 2 of 3 CGMES grids exceed $1\%$.  \textbf{SimBench
is the extreme case}: with \texttt{net.gen} empty and all $3{,}716$\,MW in
\texttt{sgen}, sweeping the load scale from $0.60$ to $1.40$ moved its total
injection from $3{,}669.9$ to $3{,}588.2$\,MW --- a $2.2\%$ swing where the
preset intends $\pm40\%$.  Its $36{,}000$ v1 rows are close to duplicates of one
another, which is also why it converged $100\%$ of the time.

\paragraph{Semantics.}  Static generation now follows the same global scale as
\texttt{net.gen}, with $q$ scaled alongside $p$ so the unit's power factor is
preserved, and receives the same per-unit dispatch jitter with $p$ and $q$ drawn
independently (matching how loads are treated).  Unlike $p$, $q$ is signed and
is not clamped: an sgen may legitimately absorb reactive power.

\paragraph{Verification.}  Comparing injection summed over sgen buses across
rows, v1 versus v2:

\begin{center}
\begin{tabular}{lrr}
\toprule
\textbf{Grid} & \textbf{v1 spread} & \textbf{v2 spread} \\
\midrule
case89pegase    & \dbad{$0.0\%$} & \dok{$78.7\%$} \\
case2848rte     & $9.7\%$        & \dok{$79.7\%$} \\
case6470rte     & $15.2\%$       & \dok{$55.1\%$} \\
case24\_ieee\_rts & $30.9\%$     & \dok{$86.2\%$} \\
\bottomrule
\end{tabular}
\end{center}

\noindent
case89pegase is the cleanest demonstration: its sgen buses carry no load, so in
v1 their injection was constant to the digit ($2430.8$--$2430.8$\,MW).  The v2
spread of $\sim\!80\%$ is what a $[0.60, 1.40]$ scale implies.

\paragraph{Effect on yield.}  Measured at $n{=}60$ per grid, perturbing sgen
moves convergence by less than the sampling noise on every PPC grid
($\pm5\%$, mostly $0\%$).  SimBench is the exception, falling from $100\%$ to
$85$--$95\%$ --- which is not a regression but the grid finally having an
operating point that moves.

\subsection{The PV-bus reactive entry now matches \texttt{u\_start}}
\label{sec:pvq}

At a PV bus, $Q$ is not a specified quantity, and what \texttt{S\_start} carried
in v1 was not a placeholder either.  The pipeline compiles the PPC with

\begin{center}
\texttt{pp.runpp(net, init="flat", max\_iteration=1, tolerance\_mva=1e9)}
\end{center}

\noindent
whose enormous tolerance makes it converge immediately, after which pypower's
\texttt{pfsoln} back-computes each generator's $Q_g$ from whatever voltage the
PPC then holds --- $V_0$: flat magnitudes with the generator setpoints applied,
zero angles.  \texttt{makeSbus} folds that into \texttt{s\_multi}.  Verified to
machine precision on case118, case300 and case1354pegase:
\[
\texttt{s\_multi.imag}[\text{PV}] \;=\; \operatorname{Im}\!\left[V_0 \,
\overline{\mathbf{Y}V_0}\right]
\qquad (10^{-16}\ \text{to}\ 10^{-13}\pu).
\]

\noindent
The problem is that $V_0$ is \emph{not} \texttt{u\_start}.  With
\texttt{start\_mode=dc\_compile}, \texttt{u\_start} carries DC angles, so the
stored reactive entry describes a different operating point from the voltage
shipped beside it --- and inconsistently per grid: the correlation between
\texttt{s\_multi.imag}[PV] and $Q(\texttt{u\_start})$ is $0.911$ on case118 but
$0.568$ on case300.  A consumer of the data cannot see this.

In v2 the entry is recomputed as
\[
Q_{\text{start}}[\text{PV}] \;=\; \operatorname{Im}\!\left[u_{\text{start}} \,
\overline{\mathbf{Y}\,u_{\text{start}}}\right],
\]
so the column means the same thing on every grid: \emph{the reactive injection
this network implies at the initial voltage you were given}.  Active power at PV
buses is untouched --- $P$ there \emph{is} specified.  This uses only
\texttt{u\_start} and $Y$, never \texttt{u\_newton}, so it is not label leakage.

\paragraph{Verification.}  Rebuilding $Y$ from the stored branch columns with
the dataloader's own routine and evaluating the identity gives a maximum
relative deviation of $\sim\!10^{-14}$ across all files with PV buses.
Grids with no PV buses at all --- SimBench has 3 slack and 91 PQ, since all its
generation is \texttt{sgen} --- are unaffected by this change; their fix is
Section~\ref{sec:sgen}.


% ===================================================================
\section{Why Some Grids Retain Fewer Rows}
\label{sec:cliff}
% ===================================================================

Every grid has a sharp \emph{global load-scale threshold} above which the AC
power flow does not converge: $100\%$ below, $0\%$ above, with no gradual
falloff.  Measured by sweeping the load scale with all jitter disabled:

\begin{center}
\begin{tabular}{lccrr}
\toprule
\textbf{Grid} & \textbf{last $100\%$ scale} & \textbf{first $0\%$} &
$P(\text{scale} < \text{cliff})$ & \textbf{v1 yield} \\
\midrule
case6515rte     & $1.00$ & $1.05$ & $56.3\%$ & $55.9\%$ \\
case9241pegase  & $1.00$ & $1.05$ & $53.1\%$ & $50.8\%$ \\
case145         & $1.05$ & $1.10$ & $59.4\%$ & $21.9\%$ \\
case6495rte     & $1.10$ & $1.15$ & $65.6\%$ & $64.8\%$ \\
case6470rte     & $1.15$ & $1.20$ & $71.9\%$ & $71.0\%$ \\
case300         & $1.30$ & $1.40$ & $93.8\%$ & $77.8\%$ \\
case1354pegase  & $1.30$ & $1.40$ & $93.8\%$ & $90.6\%$ \\
case1888rte     & $>1.40$ & --- & $100\%$ & $100\%$ \\
case2869pegase  & $>1.40$ & --- & $100\%$ & $100\%$ \\
\bottomrule
\end{tabular}
\end{center}

\noindent
The one-parameter model ``yield $=$ probability the load scale lands below the
cliff'' reproduces five of the seven shortfall yields to within one percentage
point.  This is genuine infeasibility, not solver fragility.

Two grids carry a second mechanism.  \textbf{case145} sits on its own cliff, so
per-bus load jitter alone pushes samples across it: at scale $1.00$ it converges
$100\%$ of the time with no jitter but $63/30/17\%$ at
$\sigma_P = 0.05/0.10/0.15$.  \textbf{case9241pegase} is sensitive to the PV
\emph{setpoint}: at scale $1.00$, a fixed $1.0\pu$ setpoint converges $100\%$ of
the time while the preset's $U(0.95,1.05)$ gives $53\%$; per-bus load jitter is
irrelevant there.

A mechanism that was \emph{ruled out}: per-bus $Q$ jitter flipping a load's
sign.  That requires $z < -1/\sigma_Q = -6.67$, probability $1.3\times10^{-11}$;
the measured rate is $0.000$ on every grid.  Per-bus jitter also averages out in
aggregate, so large grids see \emph{less} total load variation, not more --- bus
count alone is not the driver.

\paragraph{Consequence for v2.}  Because the boundary is sharp in load scale,
the retained rows of a shortfall grid are already samples from
$U(0.60, \text{cliff})$ rather than $U(0.60, 1.40)$; oversampling and narrowing
the range give nearly the same data.  v2 oversamples so that all 31 grids keep
one identical preset, with \texttt{--runs} set to
$37{,}000 / (\text{measured yield})$.


% ===================================================================
\section{A Systematic DC-versus-AC Angle Bias}
\label{sec:dcbias}
% ===================================================================

\texttt{u\_start} and \texttt{u\_newton} agree exactly at the slack
($0.00^\circ$) but differ at every other bus by a large, grid-specific amount:
a mean of $2.05^\circ$ on case118, $18.57^\circ$ on case300 and $16.77^\circ$ on
case1888rte.  Removing that offset leaves a residual of only $1$--$2^\circ$.

\textbf{This is not a gauge artifact.}  A global phase rotation would preserve
every angle \emph{difference}; here the differences themselves change (up to
$9.87^\circ$ on individual branches).  It is the intrinsic modelling error of DC
power flow, integrated along the network.  DC assumes $|V| = 1\pu$,
$\sin\theta \approx \theta$ and $R = 0$; the lossless assumption is directly
visible in the slack injection, which is $47.7$\,MW in DC against $472.4$\,MW in
AC on case300 --- the entire $424.7$\,MW of network losses is missing.  Every
branch's angle drop is therefore slightly wrong, \emph{with a consistent sign},
and bus angles are sums of branch drops along a path from the slack.  Measured
against hop distance:

\begin{center}
\begin{tabular}{crr}
\toprule
\textbf{hops from slack} & \textbf{case300 mean $d$} & \textbf{case118 mean $d$} \\
\midrule
0 (slack) & $0.00^\circ$  & $0.00^\circ$ \\
1         & $2.74^\circ$  & $1.10^\circ$ \\
2         & $11.07^\circ$ & $1.46^\circ$ \\
4         & $14.96^\circ$ & $2.29^\circ$ \\
8         & $17.03^\circ$ & $3.91^\circ$ \\
12        & $19.93^\circ$ & --- \\
\bottomrule
\end{tabular}
\end{center}

\noindent
$\mathrm{corr}(\text{hop}, d) = 0.742$ on case300 and $0.695$ on case118.  Mean
per-branch error is only $0.442^\circ$ on case300; twelve hops of it accumulate
to $19.93^\circ$.  case300 is worst because its slack is a stub with a single
incident branch, its diameter is 17 hops, and it spans 13 voltage levels down to
$0.6$\,kV, where $R \sim X$ and the DC assumption fails hardest.

\textbf{Implication.}  The bias is physically real and cannot be normalised
away; \texttt{u\_start} and \texttt{u\_newton} are each correct states of their
own model.  But it means the model must learn a grid-specific, distance-dependent
correction between its input and its target, and the magnitude of that
correction varies with grid diameter, voltage spread and loss fraction --- which
makes it a genuine obstacle to cross-grid transfer.  It also masks smaller
effects: the start-jitter degradation in Section~\ref{sec:u0} is invisible until
this offset is removed.


% ===================================================================
\section{Generation Procedure}
% ===================================================================

All 31 files come from \texttt{main\_datagen\_multiproc\_improved.py} with
\texttt{--label\_solver pandapower\_nr}, \texttt{--start\_mode dc\_compile},
\texttt{--K 40}, \texttt{--convergence\_mode misinf},
\texttt{--drop\_nonconverged}, \texttt{--no\_save\_y\_matrix},
\texttt{--perturb\_sgen} and \texttt{--pv\_q\_from\_vstart}, and \emph{without}
\texttt{--rand\_u\_start}.  The 28 PPC grids additionally pass
\texttt{--use\_force\_shunt\_when\_no\_trafo}.

Unlike v1, every grid ran as a \textbf{single-node job} rather than a chunked
array with a merge step.  Chunked runs need the chunk files and the merged file
on disk simultaneously, and the vault filesystem was at its soft quota; measured
throughput made one node sufficient for all 31.

The three CGMES grids pass \texttt{--cgmes\_path},
\texttt{--cgmes\_version 2.4.15} and \texttt{--cgmes\_ignore\_errors}.  SimBench
and ENTSO-E use \texttt{--no\_cgmes\_model\_a\_cleanup}; LVN heo1 uses
\texttt{--cgmes\_model\_a\_cleanup}, \texttt{--use\_force\_shunt\_when\_no\_trafo}
and \textbf{\texttt{--base\_sn\_mva 100}}.  The last is mandatory: every LVN
CGMES export carries \texttt{net.sn\_mva}${}=1$, at which the specified
injections reach ${\sim}162{,}000\pu$ and pandapower's NR cannot reach its
tolerance at any iteration count.

Every sample is attempted independently: a fresh perturbation is drawn, a DC
power flow provides the initial voltage, pandapower NR solves the AC problem,
and the row is written only if it converged.  \textbf{Non-converged samples are
discarded, not stored}, which is why \texttt{--runs} exceeds the retained count
wherever the grid has a feasibility cliff.


% ===================================================================
\section{Row Schema}
% ===================================================================

Each row is one complete power-flow scenario: 30 columns, ZSTD-compressed.
Array-valued fields are stored as serialised NumPy buffers (\texttt{binary}),
so a row is self-contained and the grid topology is repeated per row rather
than normalised out.

\begin{center}
\small
\begin{tabular}{lll}
\toprule
\textbf{Group} & \textbf{Columns} & \textbf{Type} \\
\midrule
Sizes    & \texttt{bus\_number}, \texttt{branch\_number} & int32 \\
Base     & \texttt{gridtype}, \texttt{U\_base}, \texttt{S\_base}, \texttt{vn\_kv} & string/double \\
Bus      & \texttt{bus\_typ}, \texttt{Y\_shunt\_bus} & binary \\
Branch   & \texttt{Branch\_f\_bus}, \texttt{Branch\_t\_bus}, \texttt{Branch\_status}, & binary \\
         & \texttt{Branch\_tau}, \texttt{Branch\_shift\_deg}, \texttt{Is\_trafo}, & \\
         & \texttt{Branch\_hv\_is\_f}, \texttt{Branch\_n} & \\
Admittance & \texttt{Branch\_y\_series\_from/\_to/\_ft}, & binary \\
           & \texttt{Branch\_y\_shunt\_from/\_to}, \texttt{Y\_Lines}, \texttt{Y\_C\_Lines} & \\
\midrule
\textbf{Input}  & \texttt{u\_start}, \texttt{S\_start} & binary \\
\textbf{Label}  & \texttt{u\_newton}, \texttt{S\_newton} & binary \\
\textbf{Diag.}  & \texttt{converged}, \texttt{final\_misinf}, \texttt{nr\_iterations} & int8/double/int32 \\
\bottomrule
\end{tabular}
\end{center}

\noindent
The learning task is $\texttt{u\_start} \mapsto \texttt{u\_newton}$.  Only the
converged \emph{endpoint} is stored, not the NR iterate trajectory.
\texttt{nr\_iterations} is $0$ throughout --- the pandapower label path does not
report an iteration count --- so it carries no information here.
\texttt{Y\_matrix} is absent by design; the admittance is reconstructed from the
branch columns at load time.

\paragraph{\texttt{S\_start} is not load-only.}  It is the signed injection
before the solve: perturbed load everywhere, plus each unit's scheduled dispatch
at controllable buses, plus --- at PV buses --- the reactive entry of
Section~\ref{sec:pvq}.  At non-controllable buses \texttt{S\_start} equals
\texttt{S\_newton} to ${\sim}10^{-6}$\,MW; at slack and PV buses they differ by
construction, because slack injection and PV reactive power are solver outputs.


% ===================================================================
\section{Contents}
% ===================================================================

Table~\ref{tab:contents} lists every file.  \emph{Attempts} is \texttt{--runs};
\emph{yield} is retained rows over attempts.

{\footnotesize
\setlength{\tabcolsep}{3.5pt}
\begin{longtable}{@{}l r r r r r r c l l@{}}
\caption{The 31 v2 \texttt{\_ppNR\_} datasets, ordered by network size.}
\label{tab:contents}\\
\toprule
\textbf{Grid} & \textbf{Buses} & \textbf{Br.} & \textbf{Rows} & \textbf{Att.} &
\textbf{Yield} & \textbf{MB} & \textbf{kV lv.} & \textbf{kV range} & \textbf{Src} \\
\midrule
\endfirsthead
\toprule
\textbf{Grid} & \textbf{Buses} & \textbf{Br.} & \textbf{Rows} & \textbf{Att.} &
\textbf{Yield} & \textbf{MB} & \textbf{kV lv.} & \textbf{kV range} & \textbf{Src} \\
\midrule
\endhead
% --- generated by collect_v2_table.py, do not hand-edit ---
case4gs                    & 4 & 4 & 37{,}000 & 37{,}000 & 100.0\% & 8.2 & 1 & 230 & PPC \\
case5                      & 5 & 6 & 37{,}000 & 37{,}000 & 100.0\% & 11.0 & 1 & 230 & PPC \\
case6ww                    & 6 & 11 & 37{,}000 & 37{,}000 & 100.0\% & 13.2 & 1 & 230 & PPC \\
case9                      & 9 & 9 & 37{,}000 & 37{,}000 & 100.0\% & 18.0 & 1 & 345 & PPC \\
case14                     & 14 & 20 & 37{,}000 & 37{,}000 & 100.0\% & 30.8 & 4 & 0.208--135 & PPC \\
case24\_ieee\_rts          & 24 & 38 & 37{,}000 & 37{,}000 & 100.0\% & 52.4 & 2 & 138--230 & PPC \\
GBreducednetwork           & 29 & 99 & 37{,}000 & 37{,}000 & 100.0\% & 65.6 & 3 & 132--400 & PPC \\
case30                     & 30 & 41 & 37{,}000 & 37{,}000 & 100.0\% & 63.5 & 1 & 135 & PPC \\
case\_ieee30               & 30 & 41 & 37{,}000 & 37{,}000 & 100.0\% & 64.2 & 4 & 1--132 & PPC \\
case33bw                   & 33 & 32 & 37{,}000 & 37{,}000 & 100.0\% & 74.0 & 1 & 12.66 & PPC \\
case39                     & 39 & 46 & 37{,}000 & 37{,}000 & 100.0\% & 82.2 & 1 & 345 & PPC \\
case57                     & 57 & 80 & 37{,}034 & 37{,}100 & 99.8\% & 122.0 & 6 & 115--500 & PPC \\
case89pegase               & 89 & 210 & 36{,}999 & 37{,}000 & 100.0\% & 175.9 & 3 & 150--380 & PPC \\
SimBench                   & 94 & 131 & 38{,}649 & 45{,}000 & 85.9\% & 188.7 & 3 & 110--380 & CGMES \\
case118                    & 118 & 186 & 37{,}000 & 37{,}000 & 100.0\% & 263.1 & 3 & 138--345 & PPC \\
case145                    & 145 & 453 & 36{,}918 & 170{,}000 & 21.7\% & 290.5 & 7 & 18--500 & PPC \\
iceland                    & 189 & 206 & 37{,}112 & 37{,}500 & 99.0\% & 361.7 & 15 & 0.22--220 & PPC \\
case\_illinois200          & 200 & 245 & 37{,}000 & 37{,}000 & 100.0\% & 424.1 & 3 & 13.8--230 & PPC \\
case300                    & 300 & 411 & 37{,}930 & 48{,}000 & 79.0\% & 659.4 & 13 & 0.6--345 & PPC \\
LVN\_heo1                  & 722 & 796 & 37{,}000 & 37{,}000 & 100.0\% & 1{,}310.7 & 8 & 3--380 & CGMES \\
case1354pegase             & 1{,}354 & 1{,}991 & 37{,}038 & 41{,}000 & 90.3\% & 2{,}785.6 & 2 & 220--380 & PPC \\
case1888rte                & 1{,}888 & 2{,}531 & 37{,}000 & 37{,}000 & 100.0\% & 3{,}943.5 & 20 & 3--380 & PPC \\
GBnetwork                  & 2{,}224 & 3{,}207 & 37{,}022 & 38{,}800 & 95.4\% & 4{,}162.6 & 8 & 6.6--400 & PPC \\
case2848rte                & 2{,}848 & 3{,}776 & 37{,}000 & 37{,}000 & 100.0\% & 5{,}906.2 & 20 & 3--380 & PPC \\
case2869pegase             & 2{,}869 & 4{,}582 & 37{,}099 & 37{,}100 & 100.0\% & 6{,}027.2 & 4 & 110--380 & PPC \\
case3120sp                 & 3{,}120 & 3{,}693 & 37{,}191 & 37{,}600 & 98.9\% & 6{,}616.5 & 6 & 16--400 & PPC \\
ENTSO\_E\_RealGridTest     & 6{,}051 & 8{,}202 & 37{,}000 & 37{,}000 & 100.0\% & 12{,}475.1 & 8 & 17--380 & CGMES \\
case6470rte                & 6{,}470 & 9{,}005 & 37{,}895 & 52{,}000 & 72.9\% & 13{,}779.4 & 19 & 3--380 & PPC \\
case6495rte                & 6{,}495 & 9{,}019 & 37{,}889 & 58{,}000 & 65.3\% & 13{,}844.3 & 21 & 3--380 & PPC \\
case6515rte                & 6{,}515 & 9{,}037 & 37{,}211 & 67{,}000 & 55.5\% & 13{,}621.2 & 21 & 3--380 & PPC \\
case9241pegase             & 9{,}241 & 16{,}049 & 38{,}508 & 75{,}000 & 51.3\% & 20{,}589.8 & 9 & 110--750 & PPC \\
\midrule
\textbf{Total (31)} & & & \textbf{1{,}153{,}495} & \textbf{1{,}410{,}100} & \textbf{81.8\%} & \textbf{108{,}030.5} & & & \\
\bottomrule
\end{longtable}
}


% ===================================================================
\section{Verification}
% ===================================================================

All 31 files were checked directly against the stored bytes, not against the
flags that were passed:

\begin{itemize}[leftmargin=*]
\item \textbf{Row count} matches the filename field and the parquet metadata;
  every grid retains at least $36{,}900$ rows.
\item \textbf{\texttt{converged} $=1$} for every row of every file.
\item \textbf{Loader columns} \texttt{u\_start}, \texttt{u\_newton},
  \texttt{S\_start}, \texttt{S\_newton} present in all 31.
\item \textbf{Clean start:} $\max\bigl||u_{\text{start}}|/V_{\text{base}} - 1\bigr|$
  over PQ buses is ${\sim}10^{-16}$.
\item \textbf{PV-bus $Q$:} $\operatorname{Im}[u_{\text{start}}\overline{\mathbf{Y}u_{\text{start}}}]$
  reproduced to ${\sim}10^{-14}$ relative, with $\mathbf{Y}$ rebuilt from the
  stored branch columns.
\item \textbf{sgen perturbation:} injection summed over sgen buses spans
  ${\sim}80\%$ of its mean, against $0.0$--$30.9\%$ in v1.
\end{itemize}


% ===================================================================
\section{Known Caveats}
% ===================================================================

\begin{itemize}[leftmargin=*]
\item \textbf{The DC-versus-AC angle bias of Section~\ref{sec:dcbias} remains},
  because it is physically real.  It is grid-specific and distance-dependent,
  and is the largest single obstacle to cross-grid transfer in this corpus.
\item \textbf{Yields are not uniform.}  Grids with a low feasibility cliff spend
  many attempts per retained row, and their retained rows sample
  $U(0.60,\text{cliff})$, not $U(0.60,1.40)$.  The load-scale marginal therefore
  differs per grid even though the \emph{attempt} distribution is identical.
\item \textbf{\texttt{final\_misinf} is written at sampling time} from the
  in-memory Y-bus against the specified injection vector, not recomputed from
  the stored file, and it is in SI units (VA) --- read it relative to
  $S_{\text{base}}$.
\item \textbf{LVN heo1 carries asymmetric branch impedance}
  (\texttt{BR\_R\_ASYM}, \texttt{BR\_X\_ASYM}), so $Y_{ft} \neq Y_{tf}$.
  Consumers building their own $Y$ must derive each direction from the from/to
  admittances scaled by the voltage-base ratio rather than reusing
  \texttt{Branch\_y\_series\_ft} for both; the naive reconstruction gives a
  $9\%$ relative residual on LVN against $10^{-9}$ when done correctly.
\end{itemize}


% ===================================================================
\section{Summary}
% ===================================================================

\begin{itemize}[leftmargin=*]
\item \textbf{31 files, $1{,}153{,}495$ rows, $108.0$\,GB}, selected by the
  \texttt{u0clean\_sgenpert\_pvqvstart} tag.  $1{,}410{,}100$ scenarios were
  attempted for an overall yield of $81.8\%$, and \emph{every} grid retains at
  least $36{,}918$ rows --- against v1, where six grids fell between $18$k and
  $33$k.  This is the first internally consistent generation of this corpus.
\item \textbf{Clean DC start} everywhere.  The removed jitter degraded
  \texttt{u\_start} by $35$--$134\%$ as a predictor of \texttt{u\_newton} while
  adding no information; a retained A/B pair on case118 and case300 lets this be
  tested rather than assumed.
\item \textbf{Static generation is perturbed}, fixing a defect that froze
  $100\%$ of SimBench's generation, $56\%$ of case24\_ieee\_rts's and $17\%$ of
  LVN heo1's across all scenarios.
\item \textbf{PV-bus $Q$ is consistent with the stored voltage}, so the column
  means the same thing on every grid.
\item \textbf{Row counts in filenames are true counts}, verified against parquet
  metadata --- they were attempt counts on the single-node half of v1.
\item \textbf{Shortfall yields are explained}: each grid has a sharp load-scale
  feasibility cliff, and a one-parameter model built from it reproduces five of
  seven v1 yields to within one percentage point.
\end{itemize}

\end{document}
